\documentclass[a4paper]{article}

\usepackage{graphicx}
\usepackage{amsmath,amssymb}
\usepackage{minted}
\usepackage{multirow}
\usepackage{float}
\usepackage{todonotes}
\usepackage{forest}
\usepackage{xstring}
\usepackage{xcolor}
\usepackage[most]{tcolorbox}
\usepackage{xparse}
\usepackage{natbib}

\usetikzlibrary{graphs,graphdrawing}
\usegdlibrary{layered}

% for external graphviz
\input{/home/donkarlo/Dropbox/repo/nd_ai_project/src/nd_ai/knoweldge_representation/language/formal/computer/programming/tex/latex/code_clips/external_graphviz.tex}

% for tags
\input{/home/donkarlo/Dropbox/repo/nd_ai_project/src/nd_ai/knoweldge_representation/language/formal/computer/programming/tex/latex/parts/control_sequence/kind/semtag/semtag.tex}

% for links
\input{/home/donkarlo/Dropbox/repo/nd_ai_project/src/nd_ai/knoweldge_representation/language/formal/computer/programming/tex/latex/code_clips/seehere.tex}

\title{Architectural design patterns}
\date{}

\begin{document}
    \maketitle
    
    
    \section{Layered Architecture}
        This architecture is useful when a system can be separated into clear responsibilities such as presentation, business logic, and data access.
        It is common in enterprise applications, CRUD systems, and web backends where maintainability and separation of concerns are important.
        It works well when each layer should depend mainly on the layer below it and when strict structural order is desired.
    
    
    \section{Client--Server}
        This architecture is useful when many users or devices need to request services or data from a central system.
        It is common in web applications, database systems, email systems, and networked applications in general.
        It works well when centralized control, authentication, storage, or coordination is needed.
    
    
    \section{Microservices}
        This architecture is useful when a large system must be split into small independently deployable services.
        It is common in large-scale web platforms, cloud-native systems, and organizations with multiple development teams.
        It works well when different parts of the system need separate scaling, deployment, or technology choices.
    
    
    \section{Monolithic Architecture}
        This architecture is useful when the system is relatively small or when simplicity of development and deployment is preferred.
        It is common in early-stage products, internal tools, and applications built by small teams.
        It works well when tight integration between components is acceptable and operational complexity should be minimized.
    
    
    \section{Event-Driven Architecture}
        This architecture is useful when components should react to events instead of relying only on direct synchronous calls.
        It is common in monitoring systems, financial systems, IoT platforms, and systems with asynchronous workflows.
        It works well when loose coupling, responsiveness, and extensibility of reactions are important.
    
    
    \section{Publish--Subscribe}
        This architecture is useful when one component should broadcast information to many others without knowing them directly.
        It is common in messaging systems, notification systems, user activity streams, and distributed platforms.
        It works well when producers and consumers must remain decoupled in time and structure.
    
    
    \section{Broker Architecture}
        This architecture is useful when communication between distributed components should be mediated by a central coordination mechanism.
        It is common in middleware systems, message-oriented systems, and service integration platforms.
        It works well when routing, discovery, protocol translation, or indirect interaction is required.
    
    
    \section{Pipes and Filters}
        This architecture is useful when data should pass through a sequence of independent processing stages.
        It is common in compilers, ETL pipelines, signal processing, and data transformation workflows.
        It works well when each stage can be developed, reused, tested, and replaced independently.
    
    
    \section{Model--View--Controller}
        This architecture is useful when user interface concerns must be separated from domain logic and input handling.
        It is common in web frameworks, desktop GUI systems, and applications with interactive views.
        It works well when the same model may support multiple views or when UI structure needs disciplined separation.
    
    
    \section{Model--View--ViewModel}
        This architecture is useful when UI components depend heavily on data binding and state synchronization.
        It is common in modern desktop and mobile frameworks that support declarative interfaces and binding mechanisms.
        It works well when view logic should be separated from domain logic while keeping the UI reactive.
    
    
    \section{Hexagonal Architecture}
        This architecture is useful when the domain core must remain independent from databases, user interfaces, and external services.
        It is common in systems that emphasize testability, long-term maintainability, and infrastructure independence.
        It works well when external components should plug into the system through explicit ports and adapters.
    
    
    \section{Clean Architecture}
        This architecture is useful when core business rules must be protected from framework and infrastructure changes.
        It is common in long-lived systems, enterprise software, and codebases that require strong dependency control.
        It works well when dependencies should point inward toward stable policies rather than outward toward tools.
    
    
    \section{Onion Architecture}
        This architecture is useful when domain logic should sit at the center and infrastructure should remain peripheral.
        It is common in domain-driven systems and backend applications with strong emphasis on model purity.
        It works well when the system should be organized in concentric layers with the domain isolated from external concerns.
    
    
    \section{Service-Oriented Architecture}
        This architecture is useful when a system is built from reusable business services that may be shared across organizational boundaries.
        It is common in large enterprises, institutional IT environments, and systems integrating many legacy applications.
        It works well when interoperability, formal service contracts, and cross-department integration are important.
    
    
    \section{Peer-to-Peer}
        This architecture is useful when nodes should communicate directly and share responsibilities without a permanent central server.
        It is common in file-sharing systems, decentralized networks, and some blockchain-related infrastructures.
        It works well when decentralization, fault tolerance, and distributed resource sharing are desired.
    
    
    \section{Master--Worker}
        This architecture is useful when a large task can be divided into smaller jobs and distributed among multiple workers.
        It is common in parallel computing, batch processing, rendering farms, and distributed computation systems.
        It works well when coordination is centralized but execution should be parallelized for throughput or speed.
    
    
    \section{Space-Based Architecture}
        This architecture is useful when a system must handle very high load without relying heavily on a single central database.
        It is common in large-scale transactional systems, real-time platforms, and systems requiring elastic scalability.
        It works well when application state and processing can be distributed across in-memory data grids.
    
    
    \section{Blackboard Architecture}
        This architecture is useful when multiple specialized components must contribute partial knowledge to solve a complex problem.
        It is common in AI systems, signal interpretation, speech understanding, and complex decision-support systems.
        It works well when no single control strategy is sufficient and problem solving must emerge from shared knowledge updates.
        \bibliographystyle{plainnat}
        \bibliography{/home/donkarlo/Dropbox/repo/nd_shared_research/bibliography/bibliography.bib}
\end{document}